\documentclass[10pt,a4paper]{amsart}
\usepackage[margin=19mm]{geometry}
\usepackage{amsmath,amssymb,mathtools}
\usepackage[scr=rsfs,cal=euler]{mathalfa}
\usepackage{xfrac}
\usepackage[unicode,hidelinks,bookmarks=false]{hyperref}
\newcommand{\ie}{i.e.\ }
\newcommand{\cf}{cf.\ }
\newcommand{\resp}{respectively\ }
\newcommand{\Subsection}{\subsection}
\providecommand{\nobreakdash}{\nobreak\hbox{-}}
\setlength{\emergencystretch}{2em}
\multlinegap0pt
\usepackage{lmodern,mathtools}
\newtheorem{theorem}{Theorem}[section]
\newtheorem{proposition}[theorem]{Proposition}
\newtheorem{lemma}[theorem]{Lemma}
\newtheorem{corollary}[theorem]{Corollary}
\newtheorem{definition}[theorem]{Definition}
\newtheorem{remark}[theorem]{Remark}
\newcommand{\T}{\mathbb T}
\newcommand{\N}{\mathbb N}
\newcommand{\Z}{\mathbb Z}
\newcommand{\C}{\mathbb C}
\newcommand{\B}{\mathcal B}
\newcommand{\D}{\mathcal D}
\newcommand{\G}{\mathcal G}
\newcommand{\I}{\mathcal I}
\newcommand{\DP}{\operatorname{DP}}
\newcommand{\Id}{\operatorname{Id}}
\newcommand{\supp}{\operatorname{supp}}
\newcommand{\dist}{\operatorname{dist}}
\newcommand{\clspan}{\overline{\operatorname{span}}}
\newcommand{\norm}[1]{\left\lVert#1\right\rVert}
\newcommand{\abs}[1]{\left\lvert#1\right\rvert}
\newcommand{\restr}[2]{#1\!\upharpoonright_{#2}}
\begin{document}
\title[Large ideals in B(L1(0,1))]{Large ideals in B(L1(0,1))}
\author{Amir Bahman Nasseri}
\date{}
\maketitle
\noindent\emph{Source and licence.} Large ideals in B(L1(0,1)). Version arXiv:2609.18348v1 (CC BY 4.0). This material is distributed under the Creative Commons Attribution 4.0 International licence, http://creativecommons.org/licenses/by/4.0/, as recorded in the arXiv OAI licence metadata for this exact version and shown on the abstract page https://arxiv.org/abs/2609.18348v1. Full rights notice: licenses/arxiv-2609.18348-CC-BY-4.0.txt.\par\smallskip
\noindent\textit{Editorial scope.} Counted unit: the original \texttt{theorem} environment \texttt{thm:norming} at source lines 437--460 together with its complete proof at lines 461--473; the delivered document keeps source lines 429--474 so that every referenced label and every citation resolves inside it. Native editable source is the paper's own concatenated TeX; the AMS wrapper is editorial formatting only. No publisher class, upstream script or unrelated artwork is redistributed.
\section{Quantitatively norming power ideals}\label{sec:norming}

The properness of a closed square does not prevent it from controlling multiplication on an ambient quotient. We isolate the relevant abstract fact. For a closed two-sided ideal $J$ of a Banach algebra $A$, put
\[
 J^{[n]}=\clspan\{x_1\cdots x_n:x_i\in J\},\qquad n\geq1,
\]
and let $B_J$ denote its closed unit ball in the inherited norm.


\begin{theorem}\label{thm:norming}
Suppose $\kappa_r,\kappa_\ell>0$ satisfy, for all $a\in A$,
\begin{equation}\label{eq:detector}
 \kappa_r\|a\|\leq\sup_{x\in B_J}\|ax\|,\qquad
 \kappa_\ell\|a\|\leq\sup_{x\in B_J}\|xa\|.
\end{equation}
Then for every $n\geq1$,
\begin{equation}\label{eq:power-norming}
 \kappa_r^n\|a\|\leq\sup_{x\in B_{J^{[n]}}}\|ax\|\leq\|a\|,
 \qquad
 \kappa_\ell^n\|a\|\leq\sup_{x\in B_{J^{[n]}}}\|xa\|\leq\|a\|.
\end{equation}
In particular, every $J^{[n]}$ is essential on both sides, and left multiplication represents $A$ as a closed subalgebra of $\B(J^{[n]})$. Right multiplication gives the analogous representation of the opposite algebra $A^{\mathrm{op}}$.

If a bounded homomorphism $\Phi:A\to B$ satisfies
\[
 \|\Phi(x)\|\geq\eta\|x\|\qquad(x\in J^{[n]})
\]
for some $\eta>0$, and $A\ne\{0\}$, then
\begin{equation}\label{eq:transfer-lower}
 \|\Phi(a)\|\geq\frac{\eta\kappa_r^n}{\|\Phi\|}\|a\|
 \qquad(a\in A).
\end{equation}
\end{theorem}

\begin{proof}
Induct on $n$. Apply the right estimate in \eqref{eq:detector} to each $ax$, with $x\in B_{J^{[n]}}$, and then take suprema. Since $xy\in B_{J^{[n+1]}}$ whenever $y\in B_J$, this gives the next right estimate. The left estimate follows by the same induction, and the upper bounds follow from submultiplicativity.

The norm of left multiplication by $a$ on $J^{[n]}$ lies between $\kappa_r^n\|a\|$ and $\|a\|$. Its representation is therefore injective and bounded below, hence has closed range. Right multiplication reverses products, which accounts for the opposite algebra. If $I$ is a nonzero right ideal and $0\ne a\in I$, the right estimate produces $x\in J^{[n]}$ with $0\ne ax\in I\cap J^{[n]}$. The left-ideal assertion is symmetric. In fact, $x$ may be chosen with $\|x\|\leq1$ and $\|ax\|\geq\kappa_r^n\|a\|-\varepsilon$ for any $\varepsilon>0$.

Finally,
\[
 \eta\kappa_r^n\|a\|
 \leq\sup_{x\in B_{J^{[n]}}}\|\Phi(ax)\|
 \leq\|\Phi(a)\|\|\Phi\|,
\]
which proves \eqref{eq:transfer-lower}.
\end{proof}

\end{document}
